\section{Electric potentials and angular spectra}\label{sec:spectral}

The radial estimates control changes of boundary data.  To obtain a
convergent normalization, we also need an equation for the spatial
dependence of an electric insertion.  The equation proved below is
approximate harmonicity with a summable power error.  Its proof uses
two slicings of the square array.  Slicing parallel to the insertion
line gives an improved estimate from the vacuum spectrum; slicing
perpendicularly gives derivative bounds which carry that estimate up
to the line.

Throughout this section the lattice spacing is one in tile
coordinates, and $0<\lambda\le\pi/3$ is fixed.  We use the plane row
representation of Input~\ref{inp:rows}.  Some spectral arguments below
adapt those of \cite[Section~4, ``Finite inner functions and a measure
on dilation orbits'' and ``The height variogram'']{GFFCompanion}.
That section is stated for $\lambda>\pi/3$.  We therefore give the
arguments in the present representation, identifying the row
identities they require.  No assertion about an angular representation
outside its stated parameter range is used as an input.

\subsection{Rows, currents, and their spectral measure}

We first specify the row coordinates and the exchange identity.  Put
$v_0=\pi/2$.  A strip angle $z$, with $|\Re z|<v_0$, corresponds to
the half-plane argument $w=e^{iz}$; conversely
$z=-i\operatorname{Log}w$, using the logarithm on $\Re w>0$.
The argument of the local matrix in a homogeneous row is
\[
 u(z)=\rho(v_0-z),\qquad
 R(u)=\begin{pmatrix}
 \sin(\lambda-u)/\sin\lambda&0&0&0\\
 0&\sin u/\sin\lambda&1&0\\
 0&1&\sin u/\sin\lambda&0\\
 0&0&0&\sin(\lambda-u)/\sin\lambda
 \end{pmatrix}.
\]
On an even spatial circle of size $N$, with quantum spin space
$(\C^2)^{\otimes N}$, define the monodromy and marked row by
\[
 \mathcal M_{N,c}(z)=R_{c0}(u(z))\cdots R_{c,N-1}(u(z)),\qquad
 T_N^D(z)=\operatorname{tr}_c\bigl(D_c\mathcal M_{N,c}(z)\bigr).
\]
The index $c$ denotes a two-dimensional auxiliary space.  Let
$\Omega_N$ be the positive unit Perron vector of the untwisted
rows on the zero-total-spin sector, and let $\Lambda_N(z)$ be
its eigenvalue for real strip angles, analytically continued
using that common eigenvector.  The marked rows themselves act
on the full quantum spin space; in particular a local matrix
$O$ need not preserve the zero-total-spin sector.
For real strip angles, divide by $\Lambda_N(z)$ and take plane
word limits.  Denote the resulting row and its bounded holomorphic
continuation by $\widehat A^D(z)$.  We use
$A^D(w)=\widehat A^D(-i\operatorname{Log}w)$ below.  A finite row
word means a vector
\[
 \widehat A^{D_r}(z_r)\cdots\widehat A^{D_1}(z_1)\Omega,
 \qquad r<\infty,
\]
initially with real $z_j\in(-v_0,v_0)$; the empty word is the
vacuum $\Omega$.

The local Yang--Baxter identity, and its row consequence, are
\begin{align}
 R_{12}(u-u')R_{13}(u)R_{23}(u')
 &=R_{23}(u')R_{13}(u)R_{12}(u-u'),\notag\\
 \mathcal R_{12}(z,z')\mathcal M_{N,1}(z)\mathcal M_{N,2}(z')
 &=\mathcal M_{N,2}(z')\mathcal M_{N,1}(z)
       \mathcal R_{12}(z,z'),\label{eq:row-exchange}\\
 \mathcal R_{12}(z,z')&=R_{12}\bigl(\rho(z'-z)\bigr).\notag
\end{align}
These are the finite identities in the Row Algebra and the
Infinite-Plane State subsection of \cite{FKCompanion}.
In particular, if $\mathcal R$ is invertible and
\[
 \mathcal R(D\otimes E)\mathcal R^{-1}
 =\sum_\ell D_\ell\otimes E_\ell,
\]
then auxiliary tracing, followed by the plane limit, gives the
marked exchange rule
\begin{equation}\label{eq:marked-row-exchange}
 \widehat A^D(z)\widehat A^E(z')
 =\sum_\ell\widehat A^{E_\ell}(z')\widehat A^{D_\ell}(z).
\end{equation}
There are at most sixteen terms after expansion in tensor matrix
units.  The normalization cancels on the two sides.  Analyticity
extends this identity wherever its intertwiner is invertible.

Set
\[
 Z=\begin{pmatrix}1&0\\0&-1\end{pmatrix},\qquad
 X=\begin{pmatrix}0&1\\1&0\end{pmatrix},\qquad
 K(\zeta)=\operatorname{diag}(\zeta,\zeta^{-1}),
 \qquad |\zeta|=1,
\]
and
\[
 P_\zeta(w)=A^{K(\zeta)}(w),\qquad
 P(w)=A^I(w),\qquad J(w)=A^Z(w).
\]
The Hilbert space completion of the row words is denoted by
$\mathcal H$, and $\Omega$ has norm one.  The rows are bounded and
holomorphic for $\Re w>0$, with $\|A^D(w)\|\le2\|D\|$.
Their reflection rule is
\begin{equation}\label{eq:row-adjoint}
 \widehat A^D(z)^*
 =\widehat A^{X\overline D X}(-\overline z),\qquad
 A^D(w)^*=A^{X\overline D X}(\overline w).
\end{equation}
This is Equation \texttt{bulk:adjoint} of \cite{FKCompanion},
continued holomorphically from real strip angles.  In particular,
$J(w)^*=-J(\overline w)$.
The $P_\zeta$ rows with the same twist commute and are normal;
$P_\zeta(1)$ is self-adjoint.  Every $P_\zeta(w)$ is a
contraction in the closed right half-plane away from zero.  At its two boundary
arguments the untwisted rows are the spatial shift and its inverse:
\begin{equation}\label{eq:spectral-endpoints}
 P(i)=S,\qquad P(-i)=S^{-1},\qquad
 A^D(i)=SD_0,\qquad A^D(-i)=(XD^tX)_0S^{-1}.
\end{equation}
Here $D_0$ acts at one site and $S\Omega=\Omega$.  Dilation of the
argument is implemented by a strongly continuous unitary group fixing
$\Omega$.  Also $P(w)\Omega=\Omega$.

These facts follow from the Plane Kernels, Angle Generator, and
Scalar Continuation in a Fixed Band lemmas in
\cite[Section ``Angular transfer and its two leading modes'']{FKCompanion}.
We record the analytic estimates in the latter proof which will
also justify differentiation at an endpoint.  For one radius
$r_\lambda>0$, independent of all sufficiently large even $N$,
the two quotients
\begin{equation}\label{eq:finite-endpoint-division}
 T_N^D(z)T_N^I(z)^{-1},\qquad T_N^I(z)^{-1}T_N^D(z)
\end{equation}
are holomorphic and bounded by $C_\lambda\|D\|$ on the discs
$|z-v_0|<r_\lambda$ and $|z+v_0|<r_\lambda$.
This is Equation \texttt{bulk:division} and the crossing argument
in that proof.  On these endpoint discs $T_N^I(z)$ is invertible,
so $\Lambda_N(z)$ is nonzero there.  We use nonvanishing of its
complex continuation only on these discs.  In particular, the marked vacuum vectors have
uniform holomorphic bounds on these discs.

For clarity, the dense set used to extend these vectors consists
of finite words whose real angles lie in $(-v_0/3,v_0/3)$.
It is dense because orthogonality to such words extends, one
argument at a time, by holomorphy.  Exchanging an endpoint row
through any one such word uses
Equation~\eqref{eq:marked-row-exchange}.  Its finitely many
intertwiners stay invertible on a common disc about that endpoint:
the word angles are separated from the endpoint by the middle
third of the strip.  The continuation radius is independent of
the word, although its vector bound may depend on the word.
This is the dense-set continuation step in the proof of the same
Scalar Continuation lemma.  The denominators of those intertwiners
do not depend on $D$.  Thus the scalar continuation discs, and
the constants in the low-band arguments below, are uniform for
the unit diagonal marks $K(\zeta)$.

We will need the following concrete consequence for a fixed local
spin matrix $O$.  For the finite vacuum $\Omega_N$ just specified,
there are
some $r_O>0$ and $C_O<\infty$,
\begin{equation}\label{eq:local-endpoint-bound}
 \sup_N\sup_{|z-v_0|<r_O}
 \left\|\frac{T_N^I(z)}{\Lambda_N(z)}O\Omega_N\right\|
 \le C_O.
\end{equation}
Only sufficiently large circles containing the support of $O$
are included.  To obtain this bound, translate the fixed support
into an interval, expand $O$ in tensor matrix units, and express
each resulting vector by a finite product of positive endpoint
rows, using Equation~\eqref{eq:spectral-endpoints}.  Translation
commutes with the unmarked row.  Move the row at $z$ to the
vacuum through those endpoint rows.  The intertwiners are now
near $R(0)$, the invertible swap, so their coefficients are
bounded on a sufficiently small disc.  The remaining endpoint
rows are bounded shifts with local marks, and the marked vacuum
vector is bounded by Equation~\eqref{eq:finite-endpoint-division}.
This proves Equation~\eqref{eq:local-endpoint-bound} without a
bound depending on $N$.  Cauchy's formula consequently bounds
the $k$th vector Taylor coefficient by $C_O r_O^{-k}$ after
decreasing $r_O$.  Plane word convergence and these bounds pass
endpoint derivatives and their Gram pairings to the plane.

The untwisted commuting rows have spectral functions $f$ satisfying
\begin{equation}\label{eq:spectral-inner-input}
 |f(w)|\le1\quad(\Re w>0),\qquad
 f(\overline w)=\overline{f(w)},\qquad |f(it)|=1\quad(t\ne0).
\end{equation}
They continue through the imaginary axis away from zero.  Their zeros
$a_j$ lie in a fixed cone $\Re a_j\ge c_0|a_j|$; the same cone and
continuation statements hold for a fixed twist.  The resulting inner
factorization has the form
\begin{equation}\label{eq:general-spectral-factorization}
 f(w)=\gamma e^{-a w-b/w}
       \prod_j\frac{w-a_j}{w+\overline{a_j}}\,\gamma_j,
 \qquad a,b\ge0,\quad |\gamma|=|\gamma_j|=1,
\end{equation}
with the usual summability condition on the zeros.  This is precisely
the factorization established in the Inner Factorization and Low
Spectral Bands subsection of \cite{FKCompanion}.

Let $\mathcal D$ be the closed subspace generated by words whose marks
are diagonal.  It is invariant under $P,J$ and their adjoints, and
contains finite height-phase insertions, by
Equation~\eqref{eq:spectral-endpoints}.  All untwisted spectral
statements in the next proposition are on this subspace.

The proposition supplies two different estimates for the electric
potential.  Dilation decay first forces the spectral functions to
be finite products, and the current cocycle then supplies their
measure.  Its width tail makes modes containing both very small
and very large zeros rare; this will give the extra power in the
off-axis Laplacian estimate.  The degree moment has a separate
role: it makes the logarithmic height variogram integrable, giving
the growth bound needed for potential inversion in
Section~\ref{sec:electric}.

\begin{proposition}[The current spectrum]\label{prop:spectral}
Almost every spectral function on $\mathcal D$ is either $1$ or a
finite Blaschke product with
\begin{equation}\label{eq:finite-spectral-inner}
 f(0)=f(\infty)=1,
 \qquad f(\overline w)=\overline{f(w)}.
\end{equation}
There is a spectral vector $\mathsf h$, not necessarily of finite
norm, such that
\begin{equation}\label{eq:current-cocycle}
 J(w)\Omega=(1-f(w))\mathsf h.
\end{equation}
The measure $\mu$ given by the squared fiber norm of $\mathsf h$
is invariant under dilation.  On the cross-section of products whose
smallest zero has modulus one, it decomposes as
\begin{equation}\label{eq:current-disintegration}
 \mu=\dd s\,\nu(\dd F),\qquad f(w)=F(e^s w),
 \qquad \nu(1)<\infty.
\end{equation}
If $\mathcal W(F)$ is the difference between the largest and smallest
logarithmic zero moduli, and $2n_F$ is the degree, then, for some
$\eta_0>0$,
\begin{equation}\label{eq:current-width-tail}
 \nu\{\mathcal W>t\}\le C e^{-2\eta_0t},\qquad
 \int n_F\,\nu(\dd F)<\infty.
\end{equation}
The physical plane height has logarithmic variance.  More precisely,
after an axial symmetry, write a nonzero displacement as $(k,j)$ with
$k\ge1$ and $|j|\le k$.  Then
\begin{equation}\label{eq:spectral-variogram}
 \begin{split}
 \E\bigl(h(k,j)-h(0,0)\bigr)^2
 &=2\Re\int\bigl(1-f(1)^k f(\pm i)^{|j|}\bigr)\,\mu(\dd f)\\
 &=4\nu(1)\log k+O(1),
 \end{split}
\end{equation}
where the sign agrees with the direction of the second coordinate.
The error is uniform over these displacements.
\end{proposition}

\begin{proof}
We begin with the decay furnished by
Equation~\eqref{eq:marked-row-exchange}.  If $w=e^s$ and $w'$ is
fixed, then $z=-is$ and the exchange parameter in
Equation~\eqref{eq:row-exchange} is
$\rho(z'-z)=\rho z'+i\rho s$.  Put
$b(u)=\sin u/\sin\lambda$.  For $u=u_0+i\rho s$, the displayed
matrix gives
\begin{equation}\label{eq:normalized-intertwiner-limit}
 \frac{R(u)}{b(u)}
 =\operatorname{diag}(-e^{\pm i\lambda},1,1,-e^{\pm i\lambda})
       +O(e^{-\rho|s|})
 \qquad(s\longrightarrow\pm\infty).
\end{equation}
Indeed $\sin(\lambda-u)/\sin u=-e^{\pm i\lambda}
+O(e^{-2\rho|s|})$, while the off-diagonal entries become
$1/b(u)=O(e^{-\rho|s|})$.  The limit is diagonal and invertible,
and the normalized inverse is bounded for large $|s|$.  Scalar
normalization does not change auxiliary conjugation in
Equation~\eqref{eq:marked-row-exchange}.  This calculation holds
for every fixed $0<\lambda<\pi$.

Move a row $P(e^s)$ through a fixed diagonal word $v$ to the vacuum.
The limiting diagonal intertwiner commutes with the joint diagonal
mark.  Each exchange therefore differs from a commuting exchange by
$O(e^{-\rho|s|})$ in operator norm.  The number of exchanges is fixed,
all marked rows are bounded, and $P(e^s)\Omega=\Omega$.  Thus, with
some $\eta_0>0$ which we may decrease below $\rho$,
\begin{equation}\label{eq:diagonal-dilation-decay}
 \|(1-P(e^s))v\|\le C_v e^{-\eta_0|s|}
 \quad\text{outside a fixed bounded interval of }s.
\end{equation}
The endpoints are included by their strong limits.  Thus the decay
uses the explicit exchange rule and the target-range row bounds.

For a countable total set of diagonal words, the spectral integral of
$|1-f(e^s)|^2$ over $s$ is finite by
Equation~\eqref{eq:diagonal-dilation-decay}.  Therefore it is finite
for almost every spectral function on $\mathcal D$.  In
Equation~\eqref{eq:general-spectral-factorization}, a positive $a$
or $b$ would make $f(e^s)$ tend to zero at one end; hence $a=b=0$.
For a zero $a_j$ of modulus $r_j$, the identity
\[
 \left|\frac{t-a_j}{t+\overline{a_j}}\right|^2
 =1-\frac{4t\Re a_j}{|t+\overline{a_j}|^2}
\]
and the cone condition give fixed constants $c_1,c_2>0$ such that
$1-|f(e^s)|\ge c_1$ on $|s-\log r_j|\le c_2$.
Infinitely many zeros would escape towards zero or infinity and
yield infinitely many disjoint such intervals.  Thus there are only
finitely many zeros.  Integrability forces both endpoint limits to
be one.  Reality pairs nonreal zeros with their conjugates, and the
quotient of the endpoint limits is $(-1)^{\deg f}$, so the degree
is even.  A constant spectral function must be one.  This proves
Equation~\eqref{eq:finite-spectral-inner}.

Differentiate commutation of two equally twisted rows with respect
to their common twist at $\zeta=1$ and apply it to $\Omega$.
Linearity in the auxiliary mark gives
\begin{equation}\label{eq:twist-derivative-cocycle}
 (1-P(w'))J(w)\Omega=(1-P(w))J(w')\Omega.
\end{equation}
There is no current on the spectral set $f=1$.  To see this, let $Q$
be its spectral projection.  This set is invariant under dilation,
so $Q$ commutes with the dilation group; also $QS=SQ=Q$.
The endpoint formulas give
$J(-i)=-S^{-1}J(i)S^{-1}$.  Compressing and then dilating yields
\[
 QJ(-ir)Q=-QJ(ir)Q\qquad(r>0).
\]
These are the matching boundary values for the continuation of
$QJ(w)Q$ to the left half-plane by odd reflection.
The continued function is bounded on the punctured plane, and thus
extends across zero and infinity.  Liouville's theorem and oddness
make it zero.  Since $\Omega$ belongs to this spectral set, this
also shows $QJ(w)\Omega=0$.  On every other fiber $|f(1)|<1$, so
define $\mathsf h=J(1)\Omega/(1-f(1))$ there.
Equation~\eqref{eq:twist-derivative-cocycle} proves
Equation~\eqref{eq:current-cocycle}, first on a countable dense set
of arguments and then by analyticity and endpoint continuity.

Dilation covariance fixes the vacuum and transforms this quotient
without a scalar factor.  Its squared fiber norm therefore defines
a dilation-invariant measure $\mu$.  Normalize a nonconstant product
by making its smallest zero modulus one.  Translation invariance in
the remaining logarithmic coordinate gives
Equation~\eqref{eq:current-disintegration}, initially with a
sigma-finite transversal measure.  Every normalized product has a
fixed-depth, fixed-width dip at that first zero, whereas
\[
 \int |1-f(1)|^2\,\mu(\dd f)=\|J(1)\Omega\|^2<\infty.
\]
Integration over the dip proves $\nu(1)<\infty$.

Apply Equation~\eqref{eq:diagonal-dilation-decay} to
$v=J(1)\Omega$.  In the current measure this says
\[
 I(l):=\int |1-f(1)|^2|1-f(e^l)|^2\,\mu(\dd f)
 \le C e^{-2\eta_0l}\qquad(l\ge1).
\]
For an orbit of width $\mathcal W$, first put the argument $1$ in
the dip at its smallest zero, and then put $e^l$ in the dip at its
largest zero.  Both choices occupy intervals of fixed positive
length and both factors are bounded below.  Consequently
\[
 c\,\nu\{\mathcal W>t\}
 \le\int_{t-C}^{\infty}I(l)\,\dd l,
\]
with an adjustment for bounded $t$.  This proves the width tail.

We next justify the degree integrability needed for the height
variance.  This step records the local energy argument rather than
presuming an integrability property of the transversal measure.
For this local computation use the full untwisted representation
on $\mathcal H$, and let $E$ be its nonnegative spectral endpoint
generator
\[
 E(f)=-\left.\frac{\dd}{\dd s}\arg f(ie^s)\right|_{s=0}.
\]
In the general factorization its two exponential factors contribute
$a+b$, and each zero contributes $2\Re a_j/|i-a_j|^2\ge0$.
On $\mathcal D$ the exponential factors have already been excluded.
At finite
circumference, differentiating the endpoint swap row gives the local
link
\[
 h_{j,j+1}=-\rho\,\operatorname{swap}_{j,j+1}R'_{j,j+1}(0)
 =\begin{pmatrix}
 \rho\cot\lambda&0&0&0\\
 0&0&-\rho/\sin\lambda&0\\
 0&-\rho/\sin\lambda&0&0\\
 0&0&0&\rho\cot\lambda
 \end{pmatrix}_{j,j+1}.
\]
For a local spin matrix $O$, the vacuum normalization subtracts the
vacuum energy, and therefore
\begin{equation}\label{eq:local-energy-action}
 E(O\Omega)=\left[\sum_j h_{j,j+1},O\right]\Omega.
\end{equation}
Only commutator terms whose links meet the support of $O$ survive.
Equation~\eqref{eq:local-endpoint-bound} gives a disc and uniform
vector Taylor bounds for this fixed $O$.  Cauchy's formula
therefore passes the first derivative, and its Gram pairings with
fixed local vectors, to the plane, proving
Equation~\eqref{eq:local-energy-action}.  Iterated commutators stay local
and obey $\|E^kO\Omega\|\le C_O^{k+1}k!$, since each commutator
enlarges the support by at most two sites.  These analytic vectors
identify the closure of the local commutator operator with the
spectral generator on the local-spin subspace.  Thus
Equation~\eqref{eq:local-energy-action} also justifies its quadratic
forms and the following Fourier computation.

Write $p=\arg f(i)$ modulo $2\pi$, and let $\operatorname{sym}$
average a measure at $p$ and $-p$.  The spectral measure of
$Z_0\Omega=S^{-1}J(i)\Omega$ is $|1-f(i)|^2\mu$.
Set
\[
 D_0=-\langle\Omega,[Z_0,[Z_0,h_{0,1}]]\Omega\rangle.
\]
Since $[h_{0,1},Z_0+Z_1]=0$, the Fourier transform of the vacuum
double commutators has only three coefficients, at $0,1,-1$,
equal to $2D_0,-D_0,-D_0$.  By
Equation~\eqref{eq:local-energy-action}, the same transform is
twice the symmetrized energy-weighted spectral measure of
$Z_0\Omega$.  Dividing by $|1-e^{ip}|^2$ on compact subsets of
$p\ne0$ gives
\begin{equation}\label{eq:current-energy-sum-rule}
 2\operatorname{sym}\bigl(p_*(E\mu)\bigr)
 =D_0\frac{\dd p}{2\pi}\qquad(p\ne0).
\end{equation}
On a degree-$2n_F$ orbit the boundary phase is strictly decreasing
and makes $n_F$ full turns along the positive imaginary ray.  The
change of variables $-\dd p=E\,\dd s$ therefore gives
\[
 p_*(E\,\dd s\,\nu)=\left(\int n_F\,\nu(\dd F)\right)\dd p
 \qquad(p\ne0).
\]
One may first restrict to measurable sets of orbits with bounded
degree and width and to a bounded range of $s$, and then exhaust
these sets by monotone convergence.  Comparing with
Equation~\eqref{eq:current-energy-sum-rule} proves
$4\pi\int n_F\,\nu(\dd F)=D_0<\infty$.

It remains to recover the height variance.  A directed path of
physical or endpoint steps has arguments $w_j\in\{1,i,-i\}$ and
multipliers $t_j=f(w_j)$.  By
Equations~\eqref{eq:row-adjoint} and \eqref{eq:current-cocycle},
the left vacuum current is $-(1-f(w))\mathsf h^*$ and the right
vacuum current is $(1-f(w))\mathsf h$.  Thus, for $a<b$,
\begin{equation}\label{eq:two-current-coefficient}
 \left\langle\Omega,
 J(w_a)\left(\prod_{a<j<b}P(w_j)\right)J(w_b)\Omega\right\rangle
 =-\int(1-t_a)t_{a+1}\cdots t_{b-1}(1-t_b)\,\dd\mu.
\end{equation}
For any directed tile step, its transverse spin is $+1$ when the
arrow points along the left normal to that step.  The height on
the left of an arrow exceeds the height on its right by one, so
the increment along the directed step is minus its transverse
spin.  This convention applies to the physical step and to both
endpoint directions; the two minus signs cancel in the product.
Taking the real part
of Equation~\eqref{eq:two-current-coefficient} gives their physical
covariance in path order.  At an endpoint the
single-increment variance one equals
$2\Re\int(1-f(i))\,\dd\mu$, by $|f(i)|=1$ and
$\|Z_0\Omega\|=1$.  The corresponding physical-step identity
follows by moving the integration ray for
$(1-F(w))\dd w/w$ from $i\R_+$ to $\R_+$ on each orbit.
The small and large circular arcs vanish because
$F(0)=F(\infty)=1$; the integrable bounds below justify the
integration over orbits.  Now use the telescoping identity
\[
 1-\prod_{a=1}^m t_a
 =\sum_a(1-t_a)
  -\sum_{a<b}(1-t_a)t_{a+1}\cdots t_{b-1}(1-t_b).
\]
This proves the first equality in
Equation~\eqref{eq:spectral-variogram}, initially with truncated
orbit integrals.

Outside the zero range by logarithmic distance $d$, the factor
formula gives
\[
 |1-F(e^s)|+|1-F(\pm i e^s)|\le C n_F e^{-d},
\]
or the trivial uniform bound if it is smaller.  Thus the integral
outside the zero range enlarged at each end by
$\log k+\log(1+n_F)$ is bounded uniformly.  The additional
interval lengths cost at most
$C(1+\mathcal W+\log(1+n_F))$.  On an interval of radius
$\log k$ about any one logarithmic zero modulus $s_0$, the cone
condition gives
\[
 |F(e^s)|^k\le\exp\{-c k e^{-|s-s_0|}\}.
\]
Its integral on that interval is bounded uniformly in $k$.
Multiplication by $F(\pm i e^s)^{|j|}$ does not change the bound.
Consequently the real part of the orbit integral in
Equation~\eqref{eq:spectral-variogram} is
\[
 2\log k+O(1+\mathcal W+\log(1+n_F)).
\]
The error is integrable by Equation~\eqref{eq:current-width-tail}
and degree integrability.  These same majorants justify the ray
move and telescoping passage.  Axial symmetry of the plane law,
which follows from the physical plane identification in
Input~\ref{inp:rows}, supplies the asserted bound for every
displacement.
\end{proof}

\subsection{Electric insertions and the estimate away from their line}

Fix $p$ distinct tile vertices $u_1,\ldots,u_p$ on one tile axis,
in a disk of radius $Mn$, where $M$ and $p$ are fixed and $n\ge2$.
For real charges $t=(t_1,\ldots,t_p)$ with $\sum_jt_j=0$, define
\begin{equation}\label{eq:electric-characteristic}
 U_t=\exp\left(i\sum_jt_jh(u_j)\right),\qquad
 C(t;u)=\E U_t.
\end{equation}
Neutrality makes the expression independent of the additive height
constant.  Its loop expansion gives relative weight
\begin{equation}\label{eq:electric-loop-weight}
 v_I(t)=\frac{\cos(\lambda+\sum_{j\in I}t_j)}{\cos\lambda}
\end{equation}
to a loop enclosing exactly the point set $I$.  Indeed a diagonal
height string contributes $e^{it}$ when a counterclockwise loop
encloses charge $t$, while its two winding weights before the
insertion are $e^{i\lambda}$ and $e^{-i\lambda}$.  Summing their
orientations gives Equation~\eqref{eq:electric-loop-weight}.
The same identity follows for endpoint strings from
Equation~\eqref{eq:spectral-endpoints}; ice conservation permits
rerouting the strings.  These finite-word identities use the
plane-kernel identification in Input~\ref{inp:rows}.  Spin reversal
also gives $C(t;u)=C(-t;u)$.

Choose $b>0$ sufficiently small in terms of $p$ and $\lambda$, and
assume $\max_j|t_j|\le b$.  All factors in either winding
expansion are then positive.  The number of charged loops has mean
$O(1+\log n)$ by the buffered passage estimates: assign loops to
geometric scales about the marked points until scale $Mn$, and use
a bounded number of buffers at each scale.  Since
$\log v_I(t)\ge-Cb$, Jensen's inequality gives
\begin{equation}\label{eq:small-charge-lower-bound}
 C(t;u)\ge c n^{-Cb}.
\end{equation}
The constants do not require a lower bound on the separation of
distinct marked vertices.

Fix a tile vertex $z_0$ and define the complex-valued potential
\begin{equation}\label{eq:electric-potential}
 H(z)=\frac{\E[(h(z)-h(z_0))U_t]}{C(t;u)}.
\end{equation}
Changing $z_0$ adds a constant.  Its unit increments are bounded
uniformly in $n$: condition on the unoriented loops in the positive
tilt of Equation~\eqref{eq:electric-loop-weight}.  A marked spin
lies on one loop, and its orientation average is a ratio of two
orientation sums with denominator bounded away from zero for these
charges.  The numerator is uniformly bounded.  Proposition~\ref{prop:spectral}
also gives, for each fixed insertion experiment,
\begin{equation}\label{eq:potential-growth}
 |H(z)-H(z')|\le C_{t,u}\sqrt{1+\log(2+|z-z'|)}.
\end{equation}
This is Cauchy--Schwarz in the positive physical six-vertex law,
using $|U_t|=1$ and the logarithmic variance.

For a function on tile vertices, write
\[
 \Delta_2H(z)=H(z+2e_1)+H(z-2e_1)
             +H(z+2e_2)+H(z-2e_2)-4H(z).
\]
The factor two removes the signs of the slowly damped row modes.

\begin{lemma}[Off-axis Laplacian estimate]\label{lem:off-axis}
There are constants $c_*>0$ and $C<\infty$, depending on the fixed
parameter and on $p,M$, such that, if $z$ has distance $r\ge10$
from the axis containing the insertions, then
\begin{equation}\label{eq:off-axis-laplacian}
 |\Delta_2H(z)|\le C n^{Cb}r^{-2-c_*}.
\end{equation}
For each fixed experiment, unit increments in either coordinate
direction tend to zero when their base points tend to infinity
along either normal to the insertion axis.
\end{lemma}

\begin{proof}
Take rows parallel to the insertion axis, with time increasing
towards $z$.  On the insertion slice, $U_t$ is a finite product of
on-site diagonal unitaries, since all height differences can be
written along that slice.  The exterior vector therefore has norm
at most one.  The other side is the vacuum.  The $r+O(1)$ ordinary
rows between the two carry the damping $P(1)^k$, while translation
along the slice carries powers of the unitary $P(i)$.

For completeness, current placement in the stencil can be handled
before any inverse is introduced.  Start all height differences
at one location a bounded number of rows before the stencil.  A
step issued from the vacuum has spectral vector
$(1-f(w))\mathsf h$.  Telescoping a path of such steps gives the
difference of its endpoint translation multipliers.  Thus the
Laplacian stencil is, up to a common sign, the multiplier
\begin{equation}\label{eq:laplacian-fiber}
 f(1)^k f(i)^j B(f)\mathsf h,\qquad
 B(f)=f(1)^2+f(1)^{-2}+f(i)^2+f(i)^{-2}-4,
\end{equation}
where $k=r+O(1)$ and $j\in\Z$.  With the vacuum on the bra side,
the same calculation uses the conjugate multipliers.  The inverse
time powers in this notation cancel against already present
ordinary rows.  Thus it represents a finite sum of legitimate
nonnegative row products even when $f(1)$ is small.  Cauchy--Schwarz
and Equation~\eqref{eq:small-charge-lower-bound} reduce the lemma
to a bound on the norm of Equation~\eqref{eq:laplacian-fiber}.

Put $u_f=-\log|f(1)|$ and consider a small dyadic band
$u/2<u_f\le u$.  Evaluation of the factorization at $1$ and the
zero cone give the budget
\begin{equation}\label{eq:spectral-zero-budget}
 \sum_j\min(|a_j|,|a_j|^{-1})\le Cu.
\end{equation}
Each zero is either of modulus at most $Cu$ or of modulus at least
$(Cu)^{-1}$.  On the unit semicircle the logarithm of $f$ relative
to its approximate sign is $O(u)$.  Hence $B(f)=O(u^2)$ on the
whole band.  If all zeros are small, expansion of the factors gives
\[
 \log f(w)=a_f/w+O(u^2),\qquad |a_f|\le Cu,
 \qquad |w|=1,\ \Re w\ge0,
\]
with a real $a_f$; if all zeros are large, it instead gives
$\log f(w)=b_fw+O(u^2)$ with real $b_f$ and $|b_f|\le Cu$.
The logarithms are taken relative to the same sign; the squares
of $f$ in $B$ remove that sign.  The quadratic terms in $B$ cancel, because
$1^2+i^2=0$ and $1^{-2}+i^{-2}=0$.  Therefore
\begin{equation}\label{eq:pure-band-cancellation}
 |B(f)|\le Cu^3
 \quad\text{if the band has only small zeros or only large zeros}.
\end{equation}
The bounds on the remainders follow by summing the factor expansions;
Equation~\eqref{eq:spectral-zero-budget} bounds both their first
moments and the sums of their quadratic remainders.

We now quantify the cost of the remaining, mixed-zero band.  Along
one dilation orbit, a dyadic band occupies a bounded logarithmic
length on each exterior tail: there $u_f$ is comparable to the sum
of the small zero moduli, or to the sum of the reciprocal large
zero moduli, and these sums vary geometrically with dilation.
The intervening part has length at most $\mathcal W+C$.
In particular
\[
 \mu\{u/2<u_f\le u\}\le C\int(1+\mathcal W)\,\dd\nu\le C.
\]
A mixed-zero band requires
$\mathcal W\ge2\log(1/u)-C$.  By
Equation~\eqref{eq:current-width-tail}, its measure is at most
\[
 C\int_{\{\mathcal W\ge2\log(1/u)-C\}}
       (1+\mathcal W)\,\dd\nu
 \le C u^{2\eta_0},
\]
after decreasing $\eta_0$ if needed.  Combining this with the
$O(u^2)$ bound on that set and
Equation~\eqref{eq:pure-band-cancellation}, and then taking square
roots, proves
\begin{equation}\label{eq:laplacian-band-norm}
 \left(\int_{\{u/2<u_f\le u\}}|B(f)|^2\,\dd\mu\right)^{1/2}
 \le C u^{2+c_*}
\end{equation}
for some $c_*>0$.

On this band the time rows contribute $e^{-c r u}$.  Summing the
dyadic estimates in Equation~\eqref{eq:laplacian-band-norm} gives
$Cr^{-2-c_*}$.  On the complement of the small bands, expand the
stencil into its finitely many current increments before estimating.
The current measure of $\{u_f\ge u_0\}$ is finite for each $u_0>0$,
since $|1-f(1)|\ge1-e^{-u_0}$ there and
$\int|1-f(1)|^2\,\dd\mu<\infty$.  Cancel the at most two inverse
time powers in Equation~\eqref{eq:laplacian-fiber} first.  At least
$r-O(1)$ time rows remain, giving exponential damping on this
finite-measure set.  The constant spectral
function contributes no current.  This proves
Equation~\eqref{eq:off-axis-laplacian}.

The last assertion needs no power rate.  Each unit increment is
paired with a bounded exterior vector and contains arbitrarily many
ordinary rows acting on $J(1)\Omega$ or $J(\pm i)\Omega$.
On the nonconstant current spectrum $|f(1)|<1$, and the constant
spectrum has zero current.  Dominated convergence in its finite
current-vector norm therefore makes the increment tend to zero.
\end{proof}

\subsection{Twisted rows and finite differences near the axis}

The preceding proof uses a vacuum half-plane and thus loses its
damping at the insertion line.  A different slicing supplies bounds
there: an interval between insertions contains a constant diagonal
twist, and a step-two height increment is small on its low bands.

For a fixed unit twist define
\[
 J_\zeta(w)=A^{ZK(\zeta)}(w),\qquad
 H_\zeta=-\log|P_\zeta(1)|,\qquad
 Q_\zeta(u)=\1_{[0,u]}(H_\zeta).
\]
Split $Q_\zeta(u)$ further into the two signs of $P_\zeta(1)$.
For each fixed argument $w$, the joint spectral projection of
$\{f:f(w)=0\}$ is zero, by the Scalar Continuation in a Fixed
Band lemma, label \texttt{bulk:edge-band}, in \cite{FKCompanion}.
In particular this holds at $w=1$, so $H_\zeta$ is a densely
defined nonnegative self-adjoint operator.  The same zero budget as in
Equation~\eqref{eq:spectral-zero-budget}, now including the possible
exponential factors in Equation~\eqref{eq:general-spectral-factorization},
shows that on $Q_\zeta(u)$ both $P_\zeta(w)$ and its inverse are
uniformly bounded for
\begin{equation}\label{eq:twisted-annulus}
 A u<|w|<a/u,\qquad \Re w\ge0,
\end{equation}
and on the unit semicircle
\begin{equation}\label{eq:twisted-sign}
 P_\zeta(w)Q_\zeta^\sigma(u)
 =\sigma Q_\zeta^\sigma(u)+O(u),\qquad \sigma\in\{+1,-1\}.
\end{equation}
The constants $A,a$ and the errors are uniform over unit twists.

\begin{lemma}[Twisted finite differences]\label{lem:twisted-differences}
Suppose a row interval contains no insertion, has one fixed unit
twist, and leaves at least $l$ physical rows on each side of an
observation stencil.  For every fixed $k\ge1$, every $k$-fold
coordinate difference with steps of length two satisfies
\begin{equation}\label{eq:twisted-derivative-bound}
 |D_2^{(k)}H(z)|\le C_k n^{Cb}l^{-k},\qquad l\ge10k.
\end{equation}
The estimate is uniform in the twist and in the spatial position of
the stencil.  The directions of the differences may be mixed.
\end{lemma}

\begin{proof}
Write $U=P_\zeta(i)$.  The twisted endpoint formulas give
\[
 J_\zeta(i)=UZ_0,\qquad J_\zeta(-i)=-Z_0U^{-1},\qquad
 J_\zeta(-i)=-U^{-1}J_\zeta(i)U^{-1},
\]
and the analogous identity on the other ray.  This follows by
substituting $ZK(\zeta)$ in
Equation~\eqref{eq:spectral-endpoints}; $Z$ commutes with the
diagonal twist and changes sign under $X$.  Dilate the endpoint
identity before compressing.  The compressed current
$J_{\zeta,u}(w)=Q_\zeta(u)J_\zeta(w)Q_\zeta(u)$ consequently
extends to the other half of Equation~\eqref{eq:twisted-annulus} by
\begin{equation}\label{eq:twisted-current-gluing}
 J_{\zeta,u}(-w)
 =-P_\zeta(w)^{-1}J_{\zeta,u}(w)P_\zeta(w)^{-1},
\end{equation}
where the inverses are restricted to the low band.  The formulas
agree on both imaginary rays by strong endpoint continuity.
Morera's theorem applied to matrix coefficients gives a bounded
holomorphic function on the full annulus.  Its Laurent expansion
$\sum_j w^jJ_{\zeta,u,j}$ satisfies
\[
 \|J_{\zeta,u,j}\|\le C(Cu)^{|j|}.
\]
In particular its odd part on $|w|=1$ is $O(u)$.

On an equal-sign block,
Equations~\eqref{eq:twisted-sign} and
\eqref{eq:twisted-current-gluing} show that the even part is bounded
by $Cu$ times the whole block.  Absorbing this term proves that the
whole equal-sign block is $O(u)$.  A two-step height difference is,
up to its sign convention and a fixed factor, the doubled current
\[
 J_\zeta(w)P_\zeta(w)+P_\zeta(w)J_\zeta(w).
\]
For a spatial difference the precise endpoint compensation is
\[
 U^{-2}\bigl(J_\zeta(i)U+UJ_\zeta(i)\bigr)
 =U^{-1}Z_0U+Z_0.
\]
The compensating inverse is unitary and preserves the low bands.
On opposite-sign blocks its leading signs cancel, leaving $O(u)$
by boundedness of the current.  Thus the doubled current is $O(u)$
on the full low band.  This argument does not require a separate
decay bound for an undoubled opposite-sign current.

We describe why these row estimates apply to observations at any
spatial position.  In a finite row the auxiliary mark measures the
spin at the insertion seam.  Differentiate
$[P_\zeta(w),U]=0$ in $\theta$, where $\zeta=e^{i\theta}$.
Since $\partial_\theta P_\zeta(w)=iJ_\zeta(w)$ and
$\partial_\theta U=iUZ_0$, the result is the exact identity
\begin{equation}\label{eq:translated-current}
 U^{-1}J_\zeta(w)U
 =J_\zeta(w)+Z_0P_\zeta(w)-P_\zeta(w)Z_0.
\end{equation}
The local conservation law
$[Z_c+Z_0,R_{c0}]=0$ identifies the right-hand side as the row
with the measured auxiliary spin moved across tile $0$, while the
diagonal twist remains at the original seam.  Iteration gives the
probe at site $j$ by conjugation with $U^j$.  Equivalently, the
diagonal factors acquired in a translated seam move it back by ice
conservation.  A spatial height step is the on-site operator $Z$,
which commutes with these factors.  These are identities of a
finite tensor array with no wraparound; only finitely many
translations occur in every plane correlation.  They therefore
pass through the plane-kernel limit.  All compensating shifts are
unitary, so their number introduces no factor into the estimate.

On a low band, Equation~\eqref{eq:twisted-sign} gives
$P_\zeta(i)^2=I+O(u)$ and $P_\zeta(1)^2=I+O(u)$.
Accordingly each additional spatial two-step difference, acting
by translation on both sides of the probe, contributes another
factor $u$: for its translation $T$ and compressed probe $O$, use
$T^{-1}OT-O=T^{-1}[O,T-I]$.  A temporal two-step difference shifts two rows from
one side of the probe to the other and has the same property.
The inverses used in this low-band calculation are bounded there;
before compression, all expressions retain their common
nonnegative damping rows.  Induction on the number of differences
therefore gives $C_k u^k$ for a $k$-fold difference of $H$, before
division by $C(t;u)$, when both adjacent states have cutoff $u$.
For unequal cutoffs the estimate holds with their maximum, since
both projections lie in that larger low band.

The remaining rows turn these band estimates into spatial estimates.
Split each side into $[0,l^{-1}]$, dyadic bands above $l^{-1}$ up
to a small fixed threshold, and the remaining spectral projection.
A band of lower endpoint $u$ is damped by $e^{-lu}$ by its $l$
physical rows.  The low-band terms are bounded by the convergent
double dyadic sum
\[
 C_k\sum_{u,u'}\max(u,u',l^{-1})^k e^{-c l(u+u')}
 \le C_k l^{-k}.
\]
For the remaining bands, expand the finitely many differences and
keep $l-O(k)$ rows; their contribution decays exponentially.
All exterior strings are products of unit-twist rows and diagonal
unitaries.  Their norms are at most one.  Finally divide by
$C(t;u)\ge c n^{-Cb}$ to obtain
Equation~\eqref{eq:twisted-derivative-bound}.
\end{proof}

\begin{proposition}[A summable Laplacian error]\label{prop:laplacian}
For the neutral small-charge experiments of
Equation~\eqref{eq:electric-characteristic}, there are $c_1>0$ and
$C<\infty$ such that
\begin{equation}\label{eq:laplacian}
 |\Delta_2H(z)|\le C n^{Cb}
       \left(1+\min_j|z-u_j|\right)^{-2-c_1}
 \qquad(z\in\Z^2).
\end{equation}
The constants are uniform over the stated geometry and charges.
\end{proposition}

\begin{proof}
Let $r=\min_j|z-u_j|$, and first suppose $r$ is sufficiently large.
Choose $\delta>0$ so small that
$(2+c_*)(1-\delta)>2$, where $c_*$ is supplied by
Lemma~\ref{lem:off-axis}.  If the transverse distance of $z$ from
the insertion axis is at least $r^{1-\delta}$, that lemma gives
the result immediately with exponent
$(2+c_*)(1-\delta)-2>0$.

Otherwise move $z$ in the normal direction to a vertex $z'$ in
the same coordinate-parity coset whose transverse distance lies
between $r^{1-\delta}$ and $r^{1-\delta}+2$.  Its longitudinal
distance from every insertion is at least a fixed multiple of $r$.
The entire intervening stencil therefore has, in the slicing with
time along the insertion axis, an insertion-free interval with
$l\ge c r$ rows on both sides.  Each normal step of length two
changes $\Delta_2H$ by a sum of third differences, since each
central second difference is a translate of two successive
step-two differences.  By
Lemma~\ref{lem:twisted-differences}, summing at most
$C r^{1-\delta}$ steps gives
\[
 |\Delta_2H(z)-\Delta_2H(z')|
 \le C n^{Cb}r^{1-\delta}r^{-3}
 =C n^{Cb}r^{-2-\delta}.
\]
At $z'$ use Lemma~\ref{lem:off-axis}.  Taking
$c_1=\min\{\delta,(2+c_*)(1-\delta)-2\}>0$ proves the claim for
large $r$.  For bounded $r$, the uniform bound on unit increments
of $H$ proves the same estimate after enlarging $C$.
\end{proof}

The power in Equation~\eqref{eq:laplacian} makes its right-hand side
summable on each of the four coordinate-parity cosets.  Together
with the sublogarithmic growth in
Equation~\eqref{eq:potential-growth}, this is the input needed to
recover the potential from its discrete Laplacian.
